\documentclass[14pt,letterpaper]{extarticle}
\newcommand{\packetlevel}{K–1}
\newcommand{\packetid}{F03-K-v5}
% Shared layout for the Week 3 packets.  Each packet defines \packetlevel and \packetid first.
\usepackage[T1]{fontenc}
\usepackage[utf8]{inputenc}
\usepackage[default]{sourcesanspro}
\usepackage[letterpaper, left=0.5in, right=0.5in, top=0.72in, bottom=0.68in,
            headheight=18pt, headsep=0.2in, footskip=0.34in]{geometry}
\usepackage{xcolor}
\usepackage{tikz}
\usetikzlibrary{arrows.meta}
\usepackage{fancyhdr}
\usepackage{array}

\definecolor{pbgreen}{HTML}{3FAE4A}
\definecolor{pbblue}{HTML}{2F6FD6}
\definecolor{pbred}{HTML}{D8343C}
\definecolor{pbyellow}{HTML}{F5C518}
\definecolor{pbpurple}{HTML}{7D46A8}
\definecolor{pbpink}{HTML}{F07AB4}
\definecolor{pbteal}{HTML}{1EA59B}
\definecolor{pbgray}{HTML}{8C9099}
\definecolor{slotfill}{gray}{0.975}
\definecolor{framecol}{gray}{0.6}

\pagestyle{fancy}
\fancyhf{}
\fancyhead[L]{\normalsize Week 3 / Shuffle machines / \packetlevel}
\fancyfoot[L]{\small Bellingham Math Circle / Week 3 / \packetid}
\fancyfoot[R]{\small \thepage}
\renewcommand{\headrulewidth}{0pt}
\renewcommand{\footrulewidth}{0pt}

\setlength{\parindent}{0pt}
\setlength{\parskip}{0pt}
\raggedright
\pagenumbering{arabic}

\newcommand{\labelfont}{\small}
\newcommand{\cardfont}{\normalsize\bfseries}

\newcounter{prob}
\newcommand{\problem}[1]{\stepcounter{prob}\par\textbf{Problem \theprob:} #1\par}
% a diagram, centred, with space above it
\newcommand{\fig}[2][10pt]{\par\vspace{#1}{\centering\input{fig/#2}\par}}

\renewcommand{\labelfont}{\large}

\begin{document}

One turn: move every block down its own arrow. Then slide the bottom row straight up into the top row.
When you draw a machine, give every top slot one arrow going out and every bottom slot one arrow coming in.

\vspace{12pt}
\problem{Put each block in the top slot under its picture. Do turns until every block is back under its own picture, and make a tally mark for each turn.}

\fig{k-p1a}
\fig[12pt]{k-p1b}

\newpage
\problem{Do the same with four blocks on each of these machines.}

\fig{k-p2a}
\fig[18pt]{k-p2b}

\newpage
\problem{Do the same on these two machines.}

\fig{k-p3a}
\fig[14pt]{k-p3b}

\newpage
\problem{Draw arrows on each machine so the blocks get back under their pictures in the number of turns written beside it. Your partner checks each one with the blocks.}

\fig{k-p4-2}
\fig[18pt]{k-p4-3}

\newpage
\fig[0pt]{k-p4-4}

\newpage
\problem{Can you draw a machine where only one block ends up in a different slot? Try each drawing with the blocks, and if it can't be done, explain why.}

\fig{k-p5a}
\fig[24pt]{k-p5b}

\newpage
\problem{Put each block under its picture on the first machine and do one turn. Move the row of blocks to the second machine without changing its order. Draw arrows on the second machine so one more turn puts every block back under its own picture.}

\fig{k-p6a}
\fig[24pt]{k-p6b}

\newpage
\problem{Draw all the different machines with 3 slots. How many are there?}

\fig[18pt]{k-p7}

\newpage
\problem{Each partner draws a new machine with 5 slots that takes 6 turns, not the one in Problem 3. Check each other's machines with the blocks.}

\fig{k-p8a}

\newpage
\fig[0pt]{k-p8b}

\end{document}
